\documentclass[11pt]{article}
\usepackage[utf8]{inputenc}
\usepackage[T1]{fontenc}
\usepackage{microtype}
\emergencystretch=2em
\usepackage{lmodern}
\usepackage[margin=2.5cm]{geometry}
\usepackage{amsmath,amssymb}
\usepackage{booktabs}
\usepackage{tabularx}
\usepackage{enumitem}
\usepackage{xcolor}
\usepackage[hidelinks]{hyperref}
\setlist{itemsep=2pt,topsep=4pt}
\newcommand{\cod}[1]{\texttt{#1}}
\setlength{\parskip}{4pt}

\title{Cómo se armaron los simulacros\\[2pt]\large Metodología, limitaciones y qué hacer sin pauta\\[2pt]\normalsize Programación Lineal MAT-24410}
\date{Noviembre 2026}
\author{}

\begin{document}
\maketitle

\section{Qué es esto}

Todo lo que se hizo para los Parciales 2 y 3 sigue el mismo camino: leer el material del
profe, medir cuánto se parecen sus exámenes a sus tareas y clases, convertir eso en
probabilidades de que salga cada tipo de pregunta, sortear exámenes con esas probabilidades y
al final escribir y verificar ejercicios concretos para cada hueco del sorteo.

El \cod{.zip} trae el código de cada paso en \cod{src/}, numerado en el orden en que se
corre. Los pasos 1 a 4 reproducen exactamente los resultados que se usaron (las mismas
similitudes y los mismos sorteos, porque las semillas están fijas). Los pasos 5 y 6 y la
carpeta \cod{herramientas/} son versiones reescritas de los scripts de búsqueda y
verificación que usé en su momento; esos se perdieron y estos hacen lo mismo. El paso 6
vuelve a resolver 12 de las respuestas de los simulacros y las 12 coinciden.

\paragraph{Datos.} Las pautas de la Prueba 1 y la Prueba 2 (seis preguntas en total), las
Tareas 1 a 4 y las Clases 7 a 16. Eso es todo. Es poca evidencia y casi todas las
limitaciones de la sección~\ref{sec:lim} vienen de ahí.

\section{El método, paso por paso}

\subsection{Extracción (\cod{01\_extraer\_textos.py})}
Cada PDF se pasa a texto con \cod{pdftotext -layout}. Las pautas se cortan en la palabra
``SOLUCI'' para quedarnos solo con los enunciados, las tareas se parten por ``Ejercicio N'' y
las clases por bloques de ejemplo.

\subsection{¿Cuánto recicla? Similitud de textos (\cod{02\_similitud.py})}
Cada pregunta de examen se compara contra cada ejercicio de tarea y cada bloque de clase. Para
cada par se calcula
\[
s = \tfrac12\big(\cos_{\text{car}} + \cos_{\text{pal}}\big),
\]
donde $\cos_{\text{car}}$ es la similitud coseno entre vectores TF-IDF de $n$-gramas de
caracteres (3 a 5, dentro de palabra) y $\cos_{\text{pal}}$ la de $n$-gramas de palabras (1 y
2). Los caracteres atrapan fórmulas y números parecidos; las palabras atrapan redacción.

Un número suelto de similitud no dice nada, porque todos los textos de PL comparten
vocabulario (``sujeto a'', ``base'', ``óptimo''). Por eso se mide primero la similitud entre
pares de ejercicios \emph{distintos} y se toma su percentil 95 como línea base. Luego:

\begin{center}
\begin{tabular}{ll}
\toprule
$s > 0.5$ & casi copia \\
línea base p95 $< s \le 0.5$ & adaptación \\
$s \le$ p95 & nueva (al menos en redacción) \\
\bottomrule
\end{tabular}
\end{center}

\begin{center}
\small
\begin{tabular}{llcc}
\toprule
Pregunta & Lo más parecido & $s$ & Clasificación \\
\midrule
P1-Q1 & T1-Ej3 & 0.204 & nueva \\
P1-Q2 & T1-Ej10 & 0.774 & casi copia \\
P1-Q3 & T1-Ej11 & 0.319 & adaptación \\
P2-Q1 & T2-Ej6 & 0.319 & adaptación \\
P2-Q2 & T3-Ej3 & 0.547 & casi copia \\
P2-Q3 & T2-Ej5 & 0.193 & nueva \\
\bottomrule
\end{tabular}

\smallskip
Línea base entre ejercicios distintos: media 0.096, p95 0.242.
\end{center}

En cada una de las seis preguntas lo más parecido es un ejercicio de tarea, no de clase. El
texto marcó 4 de 6 por encima del umbral; al leerlas a mano se juzgó que 5 de 6 nacen de una
idea de tarea (la diferencia es justo lo que TF-IDF no ve: misma idea, otra redacción). Ese
``5 de 6'' es el dato que alimenta al modelo del Parcial 3. La mezcla 2 casi copias, 2
adaptaciones y 2 nuevas es la razón de que en los sorteos el nivel de cada pregunta se elija
con probabilidades iguales.

Para el Parcial 3 se corrió también la Tarea 4 contra las Clases 12 a 16: ningún ejercicio
pasa de 0.27 con línea base p95 de 0.397. La Tarea 4 no copia las clases, así que las
preguntas nuevas del Parcial 3 se esperan desde la tarea más que desde las láminas.

\subsection{Tipos de pregunta (arquetipos)}
No se predicen preguntas concretas sino \emph{tipos}: ``dos fases'', ``holguras
complementarias con información parcial'', ``sensibilidad con $B^{-1}$'', etc. Cada tipo trae
un tiempo estimado de resolución en minutos. El Parcial 2 usa 10 tipos (A--J) y el Parcial 3
usa 11 (A--K); la lista completa está en el diccionario \cod{ARQ} de los scripts 3 y 4.

Cada ejercicio del material se etiquetó a mano con su tipo y un peso:

\begin{center}
\begin{tabular}{lc}
\toprule
Ejercicio de tarea & 2 \\
Ejemplo o ejercicio resuelto en clase & 1 \\
Teorema o resultado de clase & 0.5 \\
\bottomrule
\end{tabular}
\end{center}

La tarea pesa el doble porque es lo que el profe eligió que hicieras y porque de ahí salieron
sus exámenes. Si un ejercicio mezcla tipos se reparte (por ejemplo, Ej8 de la Tarea 4 cuenta
0.6 como sensibilidad leyendo el tableau y 0.4 como problema en contexto). Sumando se obtiene
un vector de pesos por tipo; normalizado es $\hat m$.

\subsection{Modelo del Parcial 2 (\cod{03\_modelo\_parcial2.py})}
Un examen simulado se genera así:
\begin{enumerate}
\item Se elige el escenario. Con probabilidad $\pi$ es ``molde'' (el profe repite lo que ya
preguntó) y con $1-\pi$ es ``rotación'' (evita lo que ya salió en la Prueba 2).
\item Se sortean las probabilidades de cada tipo, $\theta \sim \text{Dirichlet}(\alpha)$, con
\[
\alpha_{\text{molde}} = \beta\,\hat m + \gamma\, n, \qquad
\alpha_{\text{rot}} = \beta\,\hat m \cdot \rho^{\,\mathbf 1[\text{salió en P2}]} + 0.3\,\gamma\, n\cdot\mathbf 1[\text{no salió en P2}].
\]
$n$ cuenta lo que ya apareció en exámenes (la Prueba 1 cuenta 0.5 porque tenía otro temario),
$\beta$ dice cuánto le creemos al material, $\gamma$ cuánto a los exámenes, $\rho<1$ castiga
lo ya preguntado.
\item Número de preguntas $N\in\{2,3,4\}$ con probabilidades $0.2,\ 0.6,\ 0.2$ (los dos
exámenes tuvieron 3).
\item Se eligen $N$ tipos distintos sin reemplazo, cada uno con probabilidad proporcional a
$\theta$ entre los que quedan (muestreo de Plackett--Luce).
\item Si la suma de tiempos no cae entre 60 y 95 minutos el examen se descarta: no cabría en
hora y media o sería trivial.
\end{enumerate}
Con $20\,000$ exámenes aceptados se estima la probabilidad de que aparezca cada tipo. Como
$\beta,\gamma,\rho,\pi$ son juicios míos y no datos, se recorrió una rejilla de 54
combinaciones ($\beta\in\{3,6,12\}$, $\gamma\in\{1,2,3\}$, $\rho\in\{0.3,0.6\}$,
$\pi\in\{0.3,0.5,0.7\}$) y se reporta el mínimo y el máximo de cada probabilidad. El caso base
es $\beta=6,\gamma=2,\rho=0.6,\pi=0.5$.

\begin{center}
\small
\begin{tabular}{llccc}
\toprule
& Tipo & Base & Mín & Máx \\
\midrule
A & Dos fases & 0.51 & 0.36 & 0.59 \\
F & Simplex completo desde holguras & 0.48 & 0.29 & 0.67 \\
C & Demostración & 0.41 & 0.23 & 0.56 \\
B & Tableau con parámetros & 0.38 & 0.22 & 0.52 \\
G & Diccionario desde una base dada & 0.33 & 0.20 & 0.49 \\
H & Soluciones básicas, TFPL & 0.33 & 0.26 & 0.48 \\
D & Optimalidad de $x^*$ con parámetro & 0.23 & 0.13 & 0.36 \\
J & Forma estándar aislada & 0.21 & 0.10 & 0.37 \\
I & Construir un ejemplo & 0.11 & 0.06 & 0.18 \\
E & Óptimo finito según parámetro & 0.11 & 0.06 & 0.17 \\
\bottomrule
\end{tabular}
\end{center}

Los rasgos sueltos (que haya un inciso de teoría corta, que den la base inicial, que alguna
pregunta sea casi copia) se estimaron con la regla de Laplace: si algo pasó $s$ veces en $n$
oportunidades, $P=(1+s)/(2+n)$, que es la media de una posterior Beta$(1+s,\,1+n-s)$. El
desenlace de cada ejercicio de cuentas (óptimo único, múltiples, no acotado, infactible,
degenerado, artificial en cero) se sorteó con los conteos del material más uno.

\subsection{Modelo del Parcial 3 (\cod{04\_modelo\_parcial3.py})}
Aquí no había ningún examen previo del tema, así que la evidencia sobre el profe se tomó de
\emph{cómo} arma exámenes, no de \emph{qué} preguntó:
\begin{itemize}
\item Hay dos vectores de pesos, uno de la Tarea 4 ($\hat t$) y otro de las Clases 12--16
($\hat c$). Se sortean $\theta_t\sim\text{Dir}(\beta\hat t+0.3)$ y
$\theta_c\sim\text{Dir}(\beta\hat c+0.3)$; el $+0.3$ evita que un tipo ausente del material
tenga probabilidad cero.
\item Cada pregunta sale de la tarea con probabilidad $\pi\sim\text{Beta}(6,2)$, que es
Beta$(1,1)$ actualizada con ``5 de 6 preguntas pasadas salieron de una tarea''.
\item Antes de llenar el examen, con probabilidad $q\sim\text{Beta}(2,2)$ se reserva un lugar
para una demostración (una de las dos pruebas tuvo una).
\item Regla de cobertura: el profe anunció ``sensibilidad, dual y primal''. Con probabilidad
$\kappa=0.8$ se descarta todo examen que no toque los dos bloques, sensibilidad
\{E, F, G, H\} y dualidad \{A, B, C, D, I, J\}.
\item Mismo $N$ y misma ventana de 60 a 95 minutos.
\end{itemize}
Rejilla: $\beta\in\{3,6,12\}$, $\pi\in\{0.6,0.75,0.9\}$ fijo, $\kappa\in\{0.6,0.8,1\}$; 27
combinaciones.

\begin{center}
\small
\begin{tabular}{llccc}
\toprule
& Tipo & Base & Mín & Máx \\
\midrule
I & Demostración de dualidad & 0.64 & 0.62 & 0.67 \\
E & Contexto + dual + precios sombra & 0.42 & 0.35 & 0.50 \\
C & Holguras complementarias & 0.33 & 0.28 & 0.40 \\
A & Formular el dual & 0.30 & 0.25 & 0.35 \\
G & Sensibilidad con $B^{-1}$ & 0.30 & 0.24 & 0.33 \\
F & Sensibilidad leyendo el tableau & 0.28 & 0.25 & 0.31 \\
B & Relación primal-dual & 0.24 & 0.22 & 0.28 \\
H & Simplex dual & 0.22 & 0.17 & 0.28 \\
D & Dual gráfico & 0.22 & 0.18 & 0.26 \\
J & Dualidad en problema de dimensión $n$ & 0.22 & 0.18 & 0.25 \\
K & Punto interior & 0.08 & 0.05 & 0.11 \\
\bottomrule
\end{tabular}
\end{center}
A nivel bloque: alguna pregunta de sensibilidad (F, G o H) 0.71; alguna de holguras o
relación primal-dual (B, C, D o J) 0.78.

\subsection{Sorteo de simulacros}
Con los hiperparámetros base se simula un examen a la vez con una semilla fija (31102026
para el Parcial 2, 9112026 para el Parcial 3). Si la combinación de tipos ya salió se descarta
y se sortea otra, para que no haya dos simulacros con la misma estructura. A cada pregunta se
le sortea su nivel (casi copia, adaptación o nueva), su desenlace o el tipo de cambio de
sensibilidad, y si lleva inciso de teoría corta o de ``formule el dual''. Los puntos se reparten
proporcionales al tiempo estimado, redondeados a medio punto. Pedir más sorteos con la misma
semilla conserva los primeros: así salieron los simulacros 4 y 5 del Parcial 3 sin tocar los
tres anteriores.

\subsection{Del sorteo al ejercicio (\cod{05\_generar\_problemas.py})}
El sorteo solo dice, por ejemplo, ``dos fases, casi copia, desenlace infactible''. El ejercicio
concreto se encuentra por búsqueda aleatoria con filtros:
\begin{enumerate}
\item Proponer un problema chico (2 o 3 restricciones, 2 o 3 variables) con enteros chicos.
\item Filtrar rápido con \cod{scipy.optimize.linprog} (HiGHS) por el desenlace pedido.
\item Confirmar con el simplex exacto en fracciones de \cod{herramientas/simplex.py}, con las
convenciones del curso: minimizar, entra el costo reducido más negativo, empates con Bland, dos
fases con artificiales solo en filas $\ge$ e $=$.
\item Quedarse solo con números de examen: denominadores de a lo más 6 (3 en sensibilidad), sin
empates salvo que el sorteo pida degeneración, número de pivotes razonable.
\end{enumerate}
Para preguntas con parámetro se trabajó al revés: se fija la base que debe ser óptima, se
escriben sus costos reducidos como funciones del parámetro y se eligen los costos para que el
conjunto donde todos son $\ge0$ tenga la forma deseada (por ejemplo $[-3,-2]\cup[1,3]$ en la
Reposición 3). Para sensibilidad, \cod{herramientas/sensibilidad.py} calcula rangos de $b_i$ y
de $c_j$ y precios sombra desde la base óptima. El texto de cada pregunta lo redacté a mano
imitando la forma de escribir del profe.

\subsection{Verificación (\cod{06\_verificar\_respuestas.py})}
Cada respuesta se resuelve tres veces: con HiGHS en punto flotante, con el simplex exacto, y
se arma el dual para comprobar dualidad fuerte ($z_{\text{primal}}=z_{\text{dual}}$). El
paramétrico se revisa barriendo el parámetro en una rejilla. Este archivo contiene soluciones,
así que no lo abras antes de hacer los simulacros.

\section{Limitaciones}\label{sec:lim}

\paragraph{Dos exámenes son casi nada.} Todo el modelo descansa en seis preguntas. Un solo
examen nuevo puede mover las probabilidades mucho más que cualquier cambio de hiperparámetros.
Las tablas no son frecuencias medidas: son una forma ordenada de decir ``esto es lo que más
pesa en el material''.

\paragraph{Los hiperparámetros son opinión.} $\beta,\gamma,\rho,\pi,\kappa$, los pesos 2/1/0.5
y el $+0.3$ de suavizado los elegí yo. La rejilla muestra cuánto cambia el resultado
\emph{dentro} de los valores que probé, no fuera de ellos.

\paragraph{El etiquetado es subjetivo.} Decidir que un ejercicio es ``holguras
complementarias'' y no ``relación primal-dual'' es criterio. Otra persona etiquetaría distinto
y las probabilidades cambiarían.

\paragraph{TF-IDF mide palabras, no ideas.} Una pregunta con la misma idea que un ejercicio de
tarea pero otra redacción y otros números sale como ``nueva''. Al revés, dos ejercicios con
redacción parecida pero distinto fondo salen como copia. Por eso el origen final se decidió
leyendo, y por eso también parte de la clasificación depende de que \cod{pdftotext} y las
expresiones regulares partan bien los PDFs.

\paragraph{Los tiempos son supuestos.} Los minutos por tipo los estimé; nadie cronometró. La
ventana de 60 a 95 minutos depende de eso.

\paragraph{La Prueba 1 era de otro temario.} Se usó con peso 0.5 para el Parcial 2, que es
otra decisión arbitraria.

\paragraph{Las preguntas se tratan casi como independientes.} El modelo solo exige que el
examen quepa en tiempo y, en el Parcial 3, que cubra los dos bloques. Un profe real diseña el
examen como un todo (por ejemplo, la misma matriz para dos preguntas), y eso no está modelado.

\paragraph{Un sorteo no es un promedio.} Cinco sorteos tienen varianza. J (dualidad en
dimensión $n$) tiene probabilidad base de 0.22 y salió en 3 de los 5 simulacros del Parcial 3;
con probabilidad 0.22 por examen, que pase en 3 o más de 5 tiene probabilidad de alrededor de
7\%. No significa que J sea más probable: para decidir qué estudiar usa las tablas, no la
frecuencia en los simulacros.

\paragraph{El generador prefiere números bonitos.} Por diseño busca denominadores chicos y
pocos pivotes. Si el profe pone fracciones feas, los simulacros te van a parecer más fáciles
de lo que es el examen en aritmética.

\paragraph{La dificultad no está calibrada.} No hay forma de comparar la dificultad de un
simulacro con la de un examen real sin que el mismo grupo resuelva los dos. Los simulacros
intentan imitar extensión y tipo, no garantizan el mismo nivel.

\paragraph{Prueba hacia atrás: una sola y con un error.} Antes de tener la Prueba 2 el análisis
puso dos fases y tableau con parámetros arriba, y ambos salieron. Pero también dije que las
demostraciones eran baja prioridad y la Prueba 2 tuvo una. Fueron 2 de 3 aciertos en un solo
examen; eso no valida el método, solo dice que no falló por completo. Ojo: los modelos actuales
ya incluyen la Prueba 2 como dato, así que no se pueden volver a ``probar'' contra ella.

\section{Qué hacer cuando no hay pauta ni examen previo}

Es la situación del Parcial 3 y la de cualquier primer parcial con un profe nuevo. La idea es
la misma: sustituir la evidencia que falta por reglas explícitas, y reemplazarlas por datos en
cuanto llegue el primer examen.

\subsection{Armar la distribución de temas}
\begin{enumerate}
\item \textbf{El material es el prior.} Haz la lista de tipos de pregunta del temario y cuenta
cuántas veces aparece cada uno en tareas (peso 2), ejemplos de clase (1) y teoremas (0.5). Si
no hay tareas, usa solo la clase; si no hay nada, reparte uniforme entre los temas anunciados.
\item \textbf{El temario anunciado es casi una restricción.} Si el profe dijo qué viene, exige
que cada examen simulado toque todos los bloques anunciados (la regla $\kappa$), y baja el peso
de lo que no anunció (en el Parcial 3, punto interior se multiplicó por 0.3).
\item \textbf{Hábitos generales del profe.} Aunque los exámenes previos sean de otro tema
dicen cómo examina: cuántas preguntas pone, si copia de la tarea, si pone demostraciones, si
mete teoría corta dentro de las cuentas. Eso se pasa con la regla de Laplace, que con poca
evidencia da probabilidades moderadas y no extremas.
\item \textbf{Sin datos, hiperparámetros moderados y rejilla ancha.} Usa $\beta$ intermedio y
reporta mínimo y máximo, no solo el caso base. Para estudiar, ordena por el \emph{mínimo}: un
tema que sale arriba de 0.25 en todas las combinaciones es una apuesta más segura que uno que
sale alto solo en el caso base.
\end{enumerate}

\subsection{Estructura por defecto}
Cuando no sabes cómo arma el examen, usa lo que tuvieron las dos pruebas: tres preguntas en
hora y media, una larga de cuentas (30 a 35 minutos), una mediana (20 a 25) y una corta de
teoría o demostración (15). Los puntos proporcionales al tiempo. Cada pregunta de cuentas con
un inciso final que pide interpretar o justificar.

\subsection{Rúbrica genérica}
Para calificarte cuando no hay pauta, reparte los puntos de cada pregunta así:

\begin{center}
\small
\begin{tabularx}{\textwidth}{lcX}
\toprule
Criterio & Peso & Qué se revisa \\
\midrule
Planteamiento & 25\% & Forma estándar correcta, variables y restricciones bien definidas,
dual bien formulado, base inicial correcta. \\
Procedimiento & 40\% & Cada pivote o cálculo justificado: quién entra y por qué, prueba de
la razón mínima, $B^{-1}$ y costos reducidos bien calculados. \\
Resultado & 20\% & Solución y valor óptimo correctos, o el certificado correcto si es
infactible o no acotado. \\
Conclusión & 15\% & Se dice explícitamente qué pasó (óptimo único, múltiples, no acotado,
infactible, degenerado) y se interpreta en el contexto si lo hay. \\
\bottomrule
\end{tabularx}
\end{center}

Reglas que conviene aplicar: un error de aritmética que se arrastra se castiga una sola vez si
lo demás es consistente; un resultado correcto sin procedimiento vale a lo más el 20\% de
``resultado''; en sensibilidad, un rango sin decir qué pasa fuera de él pierde la conclusión.
Para demostraciones: 30\% hipótesis y lo que se quiere probar bien enunciados, 50\% la idea
central (qué teorema o desigualdad se usa), 20\% rigor y cierre.

\subsection{Calibrar tiempo y puntos}
Resuelve tú mismo el simulacro con cronómetro. Si te tardas más de 1.3 veces la hora y media,
el simulacro está largo: quita un inciso o cambia una pregunta por una más corta. Si te sobran
más de 30 minutos, está corto. Ajusta los minutos de cada tipo en \cod{ARQ} con lo que
mediste; así los sorteos siguientes salen mejor calibrados.

\subsection{Otras fuentes de señal}
Exámenes de semestres anteriores del mismo profe o de otros grupos del mismo curso; lo que el
profe dijo en clase (``esto viene'', ``esto no lo voy a preguntar''); a qué temas les dedicó
más clases; qué ejercicios de tarea tienen más incisos. Todo esto entra como peso extra en el
material o como ajuste a $\kappa$.

\subsection{Actualizar cuando llegue el examen real}
\begin{enumerate}
\item Etiqueta las preguntas del examen con los mismos tipos y súmalas a $n$ (evidencia de
examen) en el script del modelo.
\item Corre \cod{02\_similitud.py} con el examen nuevo para ver si copió, adaptó o inventó, y
actualiza la probabilidad de origen ($\pi$) y la de cada rasgo.
\item Lleva un registro de qué tan bien predijo el modelo. Para cada tipo $k$ con probabilidad
predicha $p_k$ y $y_k=1$ si salió, el puntaje de Brier es
$\frac{1}{K}\sum_k (p_k-y_k)^2$. Un modelo que no sabe nada y le pone 0.27 a todo (3 de 11)
saca alrededor de 0.20; si el tuyo no le gana a eso, el material no está prediciendo nada y
conviene confiar más en los hábitos generales del profe.
\end{enumerate}

\subsection{Autocalificación}
Después de hacer un simulacro, agrega tus respuestas como casos en
\cod{06\_verificar\_respuestas.py} y córrelo: te dice si el estado, $z^*$ y $x^*$ coinciden y
si el dual da lo mismo. Para sensibilidad, \cod{herramientas/sensibilidad.py} da los rangos y
precios sombra de cualquier base que le pases. Con eso y la rúbrica de arriba puedes
calificarte sin pauta.

\end{document}
